\section{Differentially Private Federated Learning}

Federated averaging \cite{mcmahan2017communication} trains a shared model
without pooling data, but the updates themselves leak information. DP-SGD
\cite{abadi2016deep} closes that gap by clipping per-sample gradients to a norm
$C$ and adding Gaussian noise of scale $\sigma C$, giving a per-step guarantee
that composes across steps. In FL this composition runs over communication
rounds, so a client with a fixed total budget faces a scheduling decision.
Accounting is commonly expressed in zero-concentrated DP
\cite{bun2016concentrated}, under which per-round costs add. All DP mechanisms
in this work are implemented with Opacus \cite{opacus2021}.

\vspace{10pt}

The utility cost of DP is severe for dense prediction tasks. Where a classifier
must place a single label, a detector must localise objects that may occupy a
handful of pixels, and the gradient signal supporting that localisation is small
relative to the injected noise. This is the regime in which the present work
operates, and Section~\ref{sec:calibration} shows that it is also the regime in
which most configurations are uninformative: DP-training an entire detector
destroys it, while DP-training only a small head on frozen public features makes
the noise irrelevant.

\section{Adaptive Privacy Budget Allocation}

\citeauthor{ChenAdaptivePrivacyMARL} \cite{ChenAdaptivePrivacyMARL} formulate
per-round allocation as a decentralised partially observable Markov decision
process, with one agent per client, state $(t,\rho_k)$, continuous action
$\sigma\in[\sigma_{\min},\sigma_{\max}]$, and reward the change in global
accuracy. They employ DDPG \cite{lillicrap2016continuous} with parameter sharing
and a value-decomposition network \cite{sunehag2018value}, and report gains over
uniform, arithmetic and exponential baselines: $91.05\%$ versus $88.03\%$ on
MNIST and $47.75\%$ versus $47.6\%$ on CIFAR-10.

\vspace{10pt}

Two features of that setup matter here. The reported CIFAR-10 gain is
$0.15$ percentage points, roughly $0.3\%$ relative, and is published without
error bars; and the agent's state contains no model feedback, so for a fixed
budget its policy is a deterministic function of $t$, that is, a fixed schedule.
The optimality result of Section~\ref{sec:theory} therefore bounds the published
method directly rather than merely bounding hand-designed schedules.

\vspace{10pt}

A related line of work adapts the noise scale from observed training dynamics
rather than from a learned policy over the budget. \citeauthor{WangDNORL}
\cite{WangDNORL} propose DNO-RL, in which a reinforcement-learning agent
monitors gradient norms and loss variation and modulates the Gaussian noise
scale online, reporting a better privacy-utility trade-off than fixed-rate
schedules. Because its state includes model feedback, this family is not
strictly covered by our bound; Section~\ref{sec:future} returns to that
distinction, which is the one route by which a learned controller could still
escape the result proved here.

\section{Byzantine-Robust Aggregation}

FedAvg has breakdown point zero: a single client can move the mean arbitrarily.
Robust aggregators replace the mean with a statistic that tolerates outliers.
Coordinate-wise median and trimmed mean \cite{yin2018byzantine} operate
per-coordinate; Krum \cite{blanchard2017machine} selects the single update
closest to its nearest neighbours; geometric-median variants interpolate between
the two. Their guarantees are derived without differential privacy in the loop,
and Section~\ref{sec:p3-results} shows that this assumption matters: under DP
these aggregators discard the averaging that suppresses the injected noise, and
can cost more accuracy than the attacks they defend against.

\vspace{10pt}

An orthogonal family detects and excludes anomalous updates before averaging.
\citeauthor{Li_2024} \cite{Li_2024} evaluate such schemes systematically,
clipping update norms adaptively with a threshold set to the median of
historical norms and then clustering by cosine similarity; the approach removes
the need for manual threshold tuning but degrades sharply under strongly
non-IID partitions. \citeauthor{10529118} \cite{10529118} reduce updates
with UMAP before clustering, reporting complete detection of sign-flipping and
LIE attacks at substantial server-side cost. \citeauthor{ZafarGANRobustFL}
\cite{ZafarGANRobustFL} avoid the need for a root dataset by training a
conditional GAN on the server to synthesise validation data. Common to this
family is that detection quality is reported far more often than the
false-positive cost borne by honest clients, a gap the present work addresses
explicitly in Table~\ref{tab:p3}.

\vspace{10pt}

Decentralised variants remove the server entirely.
\citeauthor{FangByzantineDecentralizedFL}
\cite{FangByzantineDecentralizedFL} propose a robust local aggregation filter
over a peer-to-peer graph that tightens as training proceeds, tolerating a large
malicious fraction but assuming benign-subgraph connectivity and providing no
protection against leakage during neighbour exchange.

\section{Poisoning that Exploits DP Noise}

\citeauthor{HossainDynamicDifferentialPoisoning}
\cite{HossainDynamicDifferentialPoisoning} observe that DP's own noise gives an
attacker cover: a detector must widen its acceptance band to admit honest noisy
updates, and that width is where adversarial perturbation can hide. They
introduce $\alpha$-MPELM, an adaptive attack that tunes a poisoning degree
$\gamma$ from memorised loss history, and show that it deceives norm, accuracy
and mixed detectors more effectively than random malicious device (RMD) attacks.
Their evaluation uses small models. Section~\ref{sec:p3-envelope} shows that the
premise is dimension-dependent: the relative spread of the honest update norm
scales as $1/\sqrt{2D}$, so the hiding room the attack relies on shrinks as the
model grows, and at detector scale it is absent.

\vspace{10pt}

A parallel concern is the adaptivity of the adversary rather than its
concealment. \citeauthor{SzelagaAdaptiveAdversariesSurvey}
\cite{SzelagaAdaptiveAdversariesSurvey} show that clients which re-enter the
federation while switching attack strategies defeat defences tuned to a single
fixed threat, though their evaluation is confined to small feedforward models
and IID partitions.

\section{Federated Reinforcement Learning under Privacy Constraints}

Reinforcement learning appears in this literature in a second role, as the
learning task being federated rather than the controller of the privacy budget.
\citeauthor{ZhuoFederatedDRL} \cite{ZhuoFederatedDRL} share only noised outputs
of local Q-networks, reaching centralised policy quality without exchanging raw
observations, but requiring shared transition identifiers and demonstrating only
two-party federations. \citeauthor{ZhouConcurrentFRLIoT}
\cite{ZhouConcurrentFRLIoT} apply local DP to edge resource allocation and
observe that the injected noise can act as a regulariser, while conceding that
utility degrades as the number of noisy edge hosts grows.
\citeauthor{HafeezFDRLRoboticSurgery} \cite{HafeezFDRLRoboticSurgery} combine
DP-SGD with homomorphic encryption for robotic-assisted surgery, with fixed
privacy parameters and a cooperative threat model. These works establish that DP
and federated RL compose in practice, but none of them treats the allocation of
the budget over rounds as the object of study.

\section{Gap Addressed}

The allocation literature optimises a quantity whose optimum has not been
characterised, and evaluates candidate policies at effect sizes comparable to
the run-to-run variance of the experiments used to measure them. The poisoning
literature assumes a camouflage effect whose magnitude depends on a model
dimension that is not controlled for. This thesis addresses both: it solves the
allocation problem in closed form and tests the solution under a paired protocol
precise enough to resolve the effect sizes at stake, and it measures the honest
update envelope directly at a dimension three orders of magnitude above the
models in which the camouflage premise was established.
